\QHeading{第1题　口算。}
\par\noindent


\begin{oralgrid}
\OralItem{1}{10.7-7.7}{3}
\OralItem{2}{0.47+3.98}{4.45}
\OralItem{3}{11.1-9.9}{1.2}
\OralItem{4}{27.42-18.22}{9.2}
\OralItem{5}{41.08-35}{6.08}
\OralItem{6}{8.2-8.02}{0.18}
\OralItem{7}{69.8+7.2}{77}
\OralItem{8}{43.45+9.86}{53.31}
\OralItem{9}{0.7+20.7}{21.4}
\OralItem{10}{37.8+16.6}{54.4}
\OralItem{11}{14.65-8.65}{6}
\OralItem{12}{54.14+6.36}{60.5}\end{oralgrid}
\par

\QHeading{口算思路补充（对应上表题号）}

\Answer{上面列出得数，下面按小题给出不同计算思路。计数单位法适用于小数，整数题可比较凑整、拆分或商不变等方法。}
\MethodTitle{第1小题：$10.7-7.7$}

\MethodBlock{口算思路一：按相同计数单位计算}{先看成若干个十分之一，按整数加减，最后还原成小数。}{\begin{aligned}[t]&\phantom{={}}10.7-7.7\\
&=(107-77)\div 10\\
&=30\div 10\\
&=3\end{aligned}}

\MethodBlock{口算思路二：补差法}{从减数先补到一个整数，再补到被减数，两段差相加。}{\begin{aligned}[t]&\phantom{={}}10.7-7.7\\
&=(8-7.7)+(10.7-8)\\
&=0.3+2.7\\
&=3\end{aligned}}
\MethodTitle{第2小题：$0.47+3.98$}

\MethodBlock{口算思路一：按相同计数单位计算}{先看成若干个百分之一，按整数加减，最后还原成小数。}{\begin{aligned}[t]&\phantom{={}}0.47+3.98\\
&=(47+398)\div 100\\
&=445\div 100\\
&=4.45\end{aligned}}

\MethodBlock{口算思路二：凑整再减去多加的数}{把第二个数凑成整数，再减去多加的部分。}{\begin{aligned}[t]&\phantom{={}}0.47+3.98\\
&=0.47+4-0.02\\
&=4.47-0.02\\
&=4.45\end{aligned}}
\MethodTitle{第3小题：$11.1-9.9$}

\MethodBlock{口算思路一：按相同计数单位计算}{先看成若干个十分之一，按整数加减，最后还原成小数。}{\begin{aligned}[t]&\phantom{={}}11.1-9.9\\
&=(111-99)\div 10\\
&=12\div 10\\
&=1.2\end{aligned}}

\MethodBlock{口算思路二：补差法}{从减数先补到一个整数，再补到被减数，两段差相加。}{\begin{aligned}[t]&\phantom{={}}11.1-9.9\\
&=(10-9.9)+(11.1-10)\\
&=0.1+1.1\\
&=1.2\end{aligned}}
\MethodTitle{第4小题：$27.42-18.22$}

\MethodBlock{口算思路一：按相同计数单位计算}{先看成若干个百分之一，按整数加减，最后还原成小数。}{\begin{aligned}[t]&\phantom{={}}27.42-18.22\\
&=(2742-1822)\div 100\\
&=920\div 100\\
&=9.2\end{aligned}}

\MethodBlock{口算思路二：补差法}{从减数先补到一个整数，再补到被减数，两段差相加。}{\begin{aligned}[t]&\phantom{={}}27.42-18.22\\
&=(19-18.22)+(27.42-19)\\
&=0.78+8.42\\
&=9.2\end{aligned}}
\MethodTitle{第5小题：$41.08-35$}

\MethodBlock{口算思路一：按相同计数单位计算}{先看成若干个百分之一，按整数加减，最后还原成小数。}{\begin{aligned}[t]&\phantom{={}}41.08-35\\
&=(4108-3500)\div 100\\
&=608\div 100\\
&=6.08\end{aligned}}

\MethodBlock{口算思路二：补差法}{从减数先补到一个整数，再补到被减数，两段差相加。}{\begin{aligned}[t]&\phantom{={}}41.08-35\\
&=(40-35)+(41.08-40)\\
&=5+1.08\\
&=6.08\end{aligned}}
\MethodTitle{第6小题：$8.2-8.02$}

\MethodBlock{口算思路一：按相同计数单位计算}{先看成若干个百分之一，按整数加减，最后还原成小数。}{\begin{aligned}[t]&\phantom{={}}8.2-8.02\\
&=(820-802)\div 100\\
&=18\div 100\\
&=0.18\end{aligned}}
\MethodTitle{第7小题：$69.8+7.2$}

\MethodBlock{口算思路一：按相同计数单位计算}{先看成若干个十分之一，按整数加减，最后还原成小数。}{\begin{aligned}[t]&\phantom{={}}69.8+7.2\\
&=(698+72)\div 10\\
&=770\div 10\\
&=77\end{aligned}}
\MethodTitle{第8小题：$43.45+9.86$}

\MethodBlock{口算思路一：按相同计数单位计算}{先看成若干个百分之一，按整数加减，最后还原成小数。}{\begin{aligned}[t]&\phantom{={}}43.45+9.86\\
&=(4345+986)\div 100\\
&=5331\div 100\\
&=53.31\end{aligned}}

\MethodBlock{口算思路二：凑整再减去多加的数}{把第二个数凑成整数，再减去多加的部分。}{\begin{aligned}[t]&\phantom{={}}43.45+9.86\\
&=43.45+10-0.14\\
&=53.45-0.14\\
&=53.31\end{aligned}}
\MethodTitle{第9小题：$0.7+20.7$}

\MethodBlock{口算思路一：按相同计数单位计算}{先看成若干个十分之一，按整数加减，最后还原成小数。}{\begin{aligned}[t]&\phantom{={}}0.7+20.7\\
&=(7+207)\div 10\\
&=214\div 10\\
&=21.4\end{aligned}}

\MethodBlock{口算思路二：凑整再减去多加的数}{把第二个数凑成整数，再减去多加的部分。}{\begin{aligned}[t]&\phantom{={}}0.7+20.7\\
&=0.7+21-0.3\\
&=21.7-0.3\\
&=21.4\end{aligned}}
\MethodTitle{第10小题：$37.8+16.6$}

\MethodBlock{口算思路一：按相同计数单位计算}{先看成若干个十分之一，按整数加减，最后还原成小数。}{\begin{aligned}[t]&\phantom{={}}37.8+16.6\\
&=(378+166)\div 10\\
&=544\div 10\\
&=54.4\end{aligned}}
\MethodTitle{第11小题：$14.65-8.65$}

\MethodBlock{口算思路一：按相同计数单位计算}{先看成若干个百分之一，按整数加减，最后还原成小数。}{\begin{aligned}[t]&\phantom{={}}14.65-8.65\\
&=(1465-865)\div 100\\
&=600\div 100\\
&=6\end{aligned}}

\MethodBlock{口算思路二：补差法}{从减数先补到一个整数，再补到被减数，两段差相加。}{\begin{aligned}[t]&\phantom{={}}14.65-8.65\\
&=(9-8.65)+(14.65-9)\\
&=0.35+5.65\\
&=6\end{aligned}}
\MethodTitle{第12小题：$54.14+6.36$}

\MethodBlock{口算思路一：按相同计数单位计算}{先看成若干个百分之一，按整数加减，最后还原成小数。}{\begin{aligned}[t]&\phantom{={}}54.14+6.36\\
&=(5414+636)\div 100\\
&=6050\div 100\\
&=60.5\end{aligned}}


\par\noindent

\QHeading{第2题　脱式计算，能简算的要简算。}

\CalcBlock[方法一]{1}{\begin{aligned}[t]&\phantom{={}}\color{black}1.7+5.6+8.3\\&\color{answerred}=1.7+8.3+5.6\\&\color{answerred}=(1.7+8.3)+5.6\\&\color{answerred}=10+5.6\\&\color{answerred}=15.6\end{aligned}}
\par
\MethodBlock{方法二：按运算顺序直接计算（对照）}{本法只作结果核对；本题要求简便计算，前面的简便解法仍保留。}{\begin{aligned}[t]&\phantom{={}}1.7+5.6+8.3\\
&=7.3+8.3\\
&=15.6\end{aligned}}


\CalcBlock[方法一]{2}{\begin{aligned}[t]&\phantom{={}}\color{black}4.76+2.8+5.24\\&\color{answerred}=4.76+5.24+2.8\\&\color{answerred}=(4.76+5.24)+2.8\\&\color{answerred}=10+2.8\\&\color{answerred}=12.8\end{aligned}}
\MethodBlock{方法二：按运算顺序直接计算（对照）}{本法只作结果核对；本题要求简便计算，前面的简便解法仍保留。}{\begin{aligned}[t]&\phantom{={}}4.76+2.8+5.24\\
&=7.56+5.24\\
&=12.8\end{aligned}}

\CalcBlock[方法一]{3}{\begin{aligned}[t]&\phantom{={}}\color{black}8.5-3.89-1.11\\&\color{answerred}=8.5-(3.89+1.11)\\&\color{answerred}=8.5-5\\&\color{answerred}=3.5\end{aligned}}
\MethodBlock{方法二：按运算顺序直接计算（对照）}{本法只作结果核对；本题要求简便计算，前面的简便解法仍保留。}{\begin{aligned}[t]&\phantom{={}}8.5-3.89-1.11\\
&=4.61-1.11\\
&=3.5\end{aligned}}

\CalcBlock[方法一]{4}{\begin{aligned}[t]&\phantom{={}}\color{black}7.2-4.9+2.8\\&\color{answerred}=7.2+2.8-4.9\\&\color{answerred}=(7.2+2.8)-4.9\\&\color{answerred}=10-4.9\\&\color{answerred}=5.1\end{aligned}}
\MethodBlock{方法二：先处理括号或乘除，再凑整再减去多加的数}{把第二个数凑成整数，再减去多加的部分。}{\begin{aligned}[t]&\phantom{={}}7.2-4.9+2.8\\
&=2.3+2.8\\
&=2.3+3-0.2\\
&=5.3-0.2\\
&=5.1\end{aligned}}

\CalcBlock[方法一]{5}{\begin{aligned}[t]&\phantom{={}}\color{black}17.3-(7.3+2.8)\\&\color{answerred}=17.3-7.3-2.8\\&\color{answerred}=10-2.8\\&\color{answerred}=7.2\end{aligned}}
\MethodBlock{方法二：按运算顺序直接计算（对照）}{本法只作结果核对；本题要求简便计算，前面的简便解法仍保留。}{\begin{aligned}[t]&\phantom{={}}17.3-(7.3+2.8)\\
&=17.3-10.1\\
&=7.2\end{aligned}}

\CalcBlock[方法一]{6}{\begin{aligned}[t]&\phantom{={}}\color{black}18.5-3.47-6.53\\&\color{answerred}=18.5-(3.47+6.53)\\&\color{answerred}=18.5-10\\&\color{answerred}=8.5\end{aligned}}
\MethodBlock{方法二：按运算顺序直接计算（对照）}{本法只作结果核对；本题要求简便计算，前面的简便解法仍保留。}{\begin{aligned}[t]&\phantom{={}}18.5-3.47-6.53\\
&=15.03-6.53\\
&=8.5\end{aligned}}

\par\noindent

\QHeading{第3题　解决问题。}

\Problem{① 某工程队修一条路，已经修了4.25千米，已经修了的长度比没修的少1.5千米。这条路长多少千米？}
\par


\MethodTitle{方法一}
\Answer{{\Formula{\begin{aligned}[t]&\phantom{={}}4.25+(4.25+1.5)\\&=4.25+5.75\\&=10\text{（千米）}\end{aligned}}}\par 答：这条路长10千米。}

\MethodTitle{方法二：用两段相同长度加上差额}
\Answer{把未修的也先看成4.25千米，两段合起来后再补上1.5千米。}
\Answer{\Formula{\begin{aligned}[t]&\phantom{={}}4.25\times 2+1.5\\
&=8.5+1.5\\
&=10\text{（千米）}\end{aligned}}}
\Answer{答：这条路长10千米。}
\Problem{② 小明的妈妈买了5千克西红柿，花了22.5元，买了同样多的黄瓜比西红柿多花5.55元。小明的妈妈买西红柿和黄瓜一共花了多少钱？}

\MethodTitle{方法一：简便计算}
\Answer{{\Formula{\begin{aligned}[t]&\phantom{={}}22.5+(22.5+5.55)\\&=(22.5+22.5)+5.55\\&=45+5.55\\&=50.55\text{（元）}\end{aligned}}}\par 答：一共花了50.55元。}

\MethodTitle{方法二：把两种蔬菜先按同样总价计算}
\Answer{两份都先按22.5元算，再加黄瓜多出的5.55元。}
\Answer{\Formula{\begin{aligned}[t]&\phantom{={}}22.5\times 2+5.55\\
&=45+5.55\\
&=50.55\text{（元）}\end{aligned}}}
\Answer{答：一共花了50.55元。}
